\documentclass[11pt]{article}
\usepackage[T1]{fontenc}
\usepackage{lmodern}
\usepackage[margin=0.95in]{geometry}
\usepackage{amsmath,amssymb}
\usepackage{microtype}
\usepackage{enumitem}
\usepackage{needspace}
\usepackage{xcolor}
\usepackage[most]{tcolorbox}
\usepackage[colorlinks=true,linkcolor=blue!50!black]{hyperref}

\newcommand{\C}{\mathbb{C}}
\newcommand{\R}{\mathbb{R}}
\newcommand{\T}{\mathbb{T}}
\newcommand{\inner}[2]{\langle #1,#2\rangle}
\newcommand{\abs}[1]{\lvert #1\rvert}
\newcommand{\norm}[1]{\lVert #1\rVert}
\newcommand{\Frob}{\mathrm{F}}
\DeclareMathOperator{\tr}{tr}
\DeclareMathOperator{\diag}{diag}
\DeclareMathOperator{\U}{U}
\newcommand{\key}[1]{\par\Needspace{5\baselineskip}\medskip\noindent{\normalsize\bfseries #1}\par\nobreak\vspace{-1pt}}
\setlist[enumerate]{leftmargin=1.5em,itemsep=3pt,topsep=3pt}

\newtcolorbox{recap}[1][]{colback=black!3,colframe=black!55,fonttitle=\bfseries,
  left=6pt,right=6pt,top=5pt,bottom=5pt,
  title={Step 1 in brief: everything from it that is used below},#1}

\title{\bfseries Step 2 of 5: the dimension collapse\\[5pt]
\large\normalfont Companion to \texttt{mub-outline.pdf}, Figure 1; preprint \S3.1--3.3 (Lemmas 3.1--3.3)\\[3pt]
\normalsize\itshape A version that assumes only school mathematics, with a summary of Step 1}
\author{}
\date{}

\begin{document}
\maketitle
\vspace{-2.4em}

\noindent\textbf{What is assumed.} Complex numbers and their moduli, matrix multiplication, the
conjugate transpose $A^*$ (transpose, then conjugate every entry), orthonormal bases, $2\times2$
determinants, and solving two equations in four unknowns. No eigenvalues, singular values or positive
definite matrices are assumed: the one $2\times2$ fact needed is proved in \S4.

\begin{recap}
\small
\textbf{The problem.} Two orthonormal bases of $\C^6$ are \emph{unbiased} if every scalar product
$\inner{b}{c}=\sum_j\overline{b_j}c_j$ between a vector of one and a vector of the other has modulus
$1/\sqrt6$. The theorem is that no four
orthonormal bases of $\C^6$ are pairwise unbiased.

\textbf{The object.} A unitary change of coordinates (one that preserves all lengths and angles) makes
the first basis the standard basis
$e_1,\dots,e_6$; the second is then the six columns of $H/\sqrt6$. Being unbiased to the first basis
says $\abs{H_{ij}}=1$ for all $i,j$, and being orthonormal says $H^*H=6I_6$. So the second basis is a
single matrix $H$, all of whose entries have modulus one and whose columns are orthogonal of length
$\sqrt6$.

\textbf{Dephasing.} Multiplying row $i$ of $H$ by a unit complex number, and column $j$ by another,
preserves both properties. Choose them so that the first row and the first column of $H$ are all $1$;
then eleven of the $36$ phases are fixed and $25$ remain, in the $5\times5$ block left after deleting
the first row and column.

\textbf{The seed.} The canonical ordering of Step 1 selects two rows and two columns of $H$; permuting
rows and columns (harmless) brings them to the front. Nine of the $25$ phases then sit in the first row
and the first column of that $5\times5$ block:
\[
a;\quad b_0,b_1,b_2,b_3;\quad c_0,c_1,c_2,c_3 ,
\]
with $a$ counted once. These nine numbers are the \emph{seed}; the other sixteen entries form the
$4\times4$ block $D'$.

\textbf{What Step 1 proved.} There is a \emph{finite} list of centres $u$, each a list of nine phases, such that every exact $H$,
after the row and column permutations of that arrangement, has its nine seed phases within Euclidean
distance $r_s=0.032$ cycles, that is $R_s=2\pi r_s=0.201$ radians, of those of some centre $u$. In one line, how: order
the six phases of two rows and of two columns canonically, cut each circle into $112$ bins, and
discard a box only when a necessary condition fails: $112^9\approx2.8\times10^{18}$ naive boxes become
about $1.6\times10^8$. A centre is a rational point chosen by the search and need not be the seed of any Hadamard matrix.
Step 1 says nothing about $D'$.

\textbf{The one inequality carried over.} If two phases differ by $\delta$ radians, the corresponding
unit complex numbers differ by at most $\delta$ (a chord is no longer than its arc). Summing the nine
squared differences gives the \emph{chord bound}
\begin{equation}\label{eq:chord}
\norm{\Delta E}_F^2+\norm{\Delta B}_F^2+\norm{\Delta C}_F^2\le R_s^2,
\end{equation}
where $E',B',C'$ are the $2+4$ blocks of $H'$ (written out in \S1), $E_0,B_0,C_0$ the corresponding
blocks of the centre, $\Delta E=E'-E_0$ and so on, and
$\norm{M}_F^2=\sum_{ij}\abs{M_{ij}}^2$ is the sum of squared moduli. Nothing else about the seed is
used below.
\end{recap}

\par\medskip\noindent\textbf{What this step receives.} One seed ball: nine phases known to within
$R_s=0.201$ radians of a centre, and nothing known about the sixteen entries of $D'$.

\par\smallskip\noindent\textbf{What it must produce.} A description of \emph{all} fillings of $D'$ that
can be covered by finitely many cells; the answer is that every completion gives a point of the
four-dimensional group $\U(2)$, at a computable distance.

\key{1.\ Where the unknown entries sit}

Write the matrix in $2+4$ block form, after the permutations of the canonical arrangement:
\[
H'=\begin{pmatrix}E'&B'\\ C'&D'\end{pmatrix},\qquad
E'=\begin{pmatrix}1&1\\ 1&a\end{pmatrix},\quad
B'=\begin{pmatrix}1&1&1&1\\[2pt] b_0&b_1&b_2&b_3\end{pmatrix},\quad
C'=\begin{pmatrix}1&c_0\\ \vdots&\vdots\\ 1&c_3\end{pmatrix},
\]
with $D'$ a full $4\times4$ block. The primes are only labels. Of the twenty entries of $E',B',C'$,
eleven are $1$ and nine are the seed; the sixteen entries of $D'$ are unknown.

The columns of $H'$ are orthogonal and have length $\sqrt6$, that is $H'^*H'=6I_6$; since
$H'/\sqrt6$ is then unitary, its rows are orthogonal as well, $H'H'^*=6I_6$. Reading the first
equation in $2+4$ blocks, the two blocks that involve $D'$ and that are used below are
\begin{equation}\label{eq:blocks}
E'^*B'+C'^*D'=0\quad(2\times4),\qquad B'^*B'+D'^*D'=6I_4 .
\end{equation}
(The remaining block involving $D'$ is the mirror image read from $H'H'^*=6I_6$ and is not used.)

\key{2.\ Two two-dimensional solution spaces}

$B'$ is $2\times4$, and its two rows $(1,1,1,1)$ and $(b_0,\dots,b_3)$ are not proportional: if they
were, all $b_j$ would equal one number $b$, and orthogonality of the first two rows of $H'$,
$1+a+4b=0$, would give $\abs{1+a}=4>2$, impossible since $\abs a=1$. So $B'v=0$ is two independent
equations in four unknowns and its solutions form a two-dimensional space $\ker B'$. The same
argument with the first two columns (whose block has rows $(1,1,1,1)$ and $(\bar c_0,\dots,\bar c_3)$,
so proportionality would again force $1+a+4c=0$) shows that the solutions of $C'^*u=0$ form a
two-dimensional space $\ker C'^*$.

Both arguments use that $H'$ is exact. A centre need not be a Hadamard matrix, so at the centre the
same two statements are \emph{checked} with room to spare rather than derived; if either check fails,
the seed is left unresolved and nothing below applies to it. The frames are therefore orthonormal
bases $u_1,u_2$ of $\ker B_0$ and $w_1,w_2$ of $\ker C_0^*$, the centre's planes, put side by side as
$4\times2$ matrices
\[
K_B=(u_1\;\;u_2),\qquad K_C=(w_1\;\;w_2),\qquad K_B^*K_B=K_C^*K_C=I_2 .
\]
The paper calls $K_B,K_C$ the \emph{frames}. They are computed once and held fixed for the whole
ball; all that is needed is solving two linear systems.

\key{3.\ Exact seeds: the block becomes a unitary matrix}

Suppose first that the seed is exactly the centre's. Let $v\in\C^4$ be nonzero with $B'v=0$. The
second equation of \eqref{eq:blocks} gives
\[
v^*D'^*D'v=6v^*v-v^*B'^*B'v=6\norm{v}^2,
\]
so $\norm{D'v}=\sqrt6\,\norm{v}$; and the first gives $C'^*D'v=-E'^*B'v=0$, so $D'v\in\ker C'^*$.
Thus

\begin{center}
\emph{$D'$ maps the plane $\ker B'$ onto the plane $\ker C'^*$ and multiplies every length by
exactly $\sqrt6$.}
\end{center}

\noindent Define $Y=K_C^*D'K_B$, a $2\times2$ matrix. If $v=K_Bc$ with $c\in\C^2$, then
$K_C^*D'v=Yc$, and since $D'v\in\ker C'^*$ and $K_C^*$ preserves lengths on that plane,
\[
\norm{Yc}=\norm{K_C^*D'K_Bc}=\norm{D'K_Bc}=\sqrt6\,\norm{c}\qquad(c\in\C^2).
\]
Hence $Y^*Y=6I_2$, that is $V=Y/\sqrt6$ satisfies $V^*V=I_2$, so $V\in\U(2)$. Sixteen unknown phases
have produced four real numbers.

\key{4.\ The price of knowing the seed only approximately}

Now the seed of $H'$ is merely within $R_s$ of the centre's, so the frames $K_B,K_C$ belong to the
centre and $B',C'$ differ from $B_0,C_0$ by $\Delta B,\Delta C$. Repeat the computation of \S3
without assuming the errors vanish. Substituting $D'^*D'=6I_4-B'^*B'$ into
$Y^*Y=K_B^*D'^*K_CK_C^*D'K_B$, and writing $P_C=I_4-K_CK_C^*$ (which deletes the part in the plane
$\ker C_0^*$), gives
\begin{equation}\label{eq:defect}
6I_2-Y^*Y=K_B^*B'^*B'K_B+K_B^*D'^*P_CD'K_B=M_1^*M_1+M_2^*M_2
\end{equation}
with $M_1=B'K_B$ and $M_2=P_CD'K_B$. A sum of matrices of the form $M^*M$ satisfies
$z^*(6I_2-Y^*Y)z\ge0$ for every $z\in\C^2$. Its trace (the sum of the two diagonal entries) is
\[
\tr(6I_2-Y^*Y)=\norm{B'K_B}_F^2+\norm{P_CD'K_B}_F^2
\le (b_1')^2+\Bigl(\frac{A_*b_1'+\sqrt6\,c_1'}{m_C}\Bigr)^2\le q,
\]
where $b_1'=\norm{\Delta B}_F$, $c_1'=\norm{\Delta C}_F$, $m_C>0$ is the least amount by which
$C_0^*$ stretches a vector orthogonal to $\ker C_0^*$ (checked at the centre), and
\begin{equation}\label{eq:q}
q=R_s^2\Bigl(1+\frac{A_*^2+6}{m_C^2}\Bigr).
\end{equation}
Here $A_*$ is a bound on $\norm{E'z}$: since every entry of $E'$ has modulus one, $A_*=2$ always
suffices, and the paper uses the smaller $\min(2,A_0+R_s)$, where $A_0$ is the largest amount by
which the centre's $E_0$ stretches a vector.

Here is the estimate, the only one in the step (Lemma 3.2). \emph{First term.} $B_0K_B=0$, so
$\norm{B'K_B}_F=\norm{\Delta BK_B}_F\le b_1'$. \emph{Second term.} Let $z=D'K_Bw$ with $\norm w=1$.
Since $C_0^*$ kills the plane $\ker C_0^*$, $C_0^*z=C_0^*P_Cz$, and by the definition of $m_C$
applied to $P_Cz$, which is orthogonal to $\ker C_0^*$, $\norm{P_Cz}\le\norm{C_0^*z}/m_C$. By
\eqref{eq:blocks}, $C_0^*D'K_B=-E'^*\Delta BK_B-\Delta C^*D'K_B$ (substitute $B'=B_0+\Delta B$ and
$C'=C_0+\Delta C$, and use $B_0K_B=0$), and the two terms on the right are at
most $A_*b_1'$ and $\sqrt6\,c_1'$, because $\norm{D'u}\le\sqrt6\norm u$ for every $u$ (as
$D'^*D'=6I_4-B'^*B'$). Squaring, summing over the two columns of $K_B$, and applying Cauchy--Schwarz
and \eqref{eq:chord} gives the bound.

\smallskip
\noindent Now the bookkeeping of \S3 pays off. One fact about $2\times2$ matrices is needed: if $A^*=A$,
then $A=W\diag(\mu_1,\mu_2)W^*$ with $W$ unitary and $\mu_j$ real. Indeed, for
$A=\begin{pmatrix}x&z\\ \bar z&y\end{pmatrix}$ with $x,y\in\R$, a nonzero $v$ with $Av=\mu v$ exists
exactly when $\det(A-\mu I)=0$, because a $2\times2$ matrix with determinant $0$ has a nonzero kernel;
that equation is $\mu^2-(x+y)\mu+(xy-\abs z^2)=0$, whose discriminant $(x-y)^2+4\abs z^2$ is
non-negative, so both roots are real. If the roots differ, the corresponding vectors are orthogonal,
since $(\mu_1-\mu_2)\inner{v_1}{v_2}=\inner{Av_1}{v_2}-\inner{v_1}{Av_2}=0$ when $A^*=A$; if they
coincide, $A=\mu I$ and any orthonormal pair will do. Applied to $A=Y^*Y$: the numbers
$\mu_j=\norm{Yw_j}^2$ are non-negative, $Y$ stretches by $\sqrt{\mu_j}$ along $w_j$, and
$\mu_1+\mu_2=\tr A$.

\smallskip
\noindent Put $\ell_j=6-\mu_j$. Then $\ell_j\ge0$ and, by \eqref{eq:defect},
$\ell_1+\ell_2=\tr(6I_2-Y^*Y)\le q<6$, so $\mu_j\ge6-q>0$. Set
$V=YW\diag(\mu_1^{-1/2},\mu_2^{-1/2})W^*$. Then $V^*V=I_2$, and $YW$ has orthogonal columns of
lengths $\sqrt{\mu_j}$, so $YW=W''\diag(\sqrt{\mu_1},\sqrt{\mu_2})$ for a unitary $W''$ and
\[
Y-\sqrt6\,V=W''\diag\bigl(\sqrt{\mu_1}-\sqrt6,\ \sqrt{\mu_2}-\sqrt6\bigr)W^*,
\qquad
\norm{Y-\sqrt6\,V}_F^2=\sum_{j=1}^2\bigl(\sqrt{\mu_j}-\sqrt6\bigr)^2 .
\]
Since $\sqrt6-\sqrt{\mu_j}=\ell_j/(\sqrt6+\sqrt{\mu_j})\le\ell_j/(\sqrt6+\sqrt{6-q})$,
\begin{equation}\label{eq:p}
\norm{Y-\sqrt6\,V}_F\ \le\ \frac{\ell_1+\ell_2}{\sqrt6+\sqrt{6-q}}\ \le\
p:=\frac{q}{\sqrt6+\sqrt{6-q}} .
\end{equation}
In the range used by the paper, $q<0.32$, hence $p<0.067$: the compressed block is not merely close
to a rigid motion, it is at a computed distance from one.

\key{5.\ Four parameters are few enough to cover}

$\U(2)$ is the set of $2\times2$ complex matrices $V$ with $V^*V=I_2$, that is, the linear maps of
$\C^2$ that preserve all lengths and angles. A $2\times2$ complex matrix is eight real numbers, and
$V^*V=I_2$ is four real equations, so four real parameters remain. They can be written out: with
$e(x)=e^{2\pi ix}$, every element of $\U(2)$ is
\[
V=e(\alpha)\begin{pmatrix}e(\beta)\cos2\pi t&e(\gamma)\sin2\pi t\\[1pt]
-e(-\gamma)\sin2\pi t&e(-\beta)\cos2\pi t\end{pmatrix},
\qquad 0\le\alpha\le\tfrac12,\quad 0\le t\le\tfrac14,\quad 0\le\beta,\gamma\le1,
\]
and a direct check shows that the two rows are orthonormal, so all these matrices are unitary.
Conversely, every $V\in\U(2)$ has this form: $\abs{\det V}=1$, so with $\alpha\in[0,\tfrac12]$ chosen
by $e(2\alpha)=\det V$ the matrix $V_0=e(-\alpha)V$ has determinant $1$, hence $V_0^{-1}=V_0^*$, and
comparing this with the $2\times2$ inverse formula gives $V_0=\begin{pmatrix}z&w\\-\bar w&\bar
z\end{pmatrix}$ with $\abs z^2+\abs w^2=1$; writing $z=e(\beta)\cos2\pi t$ and
$w=e(\gamma)\sin2\pi t$ with $0\le t\le\tfrac14$ gives the display. So the four parameters are four
angles, each in a bounded interval, and $\U(2)$ is closed and bounded. Step 3 cuts each of the four intervals
into finitely many cells; a product of four cells is a cell of $\U(2)$, and finitely many cells cover
all of it.

Counting the freedom: $25=9+16$ phases, nine in a ball and sixteen in $D'$. What the rest of the
proof keeps from $D'$ is only $Y$: every completion produces one, and $Y^*Y=6I_2$ cuts its eight real
entries down to four; many
blocks give the same $Y$, which costs nothing, since a necessary description may forget information.
So a completion is described by nine numbers in a ball and a point of a closed and bounded
four-dimensional group --- small enough to cover by finitely many cells, as a sixteen-dimensional
space would not be. Nothing here is special to dimension six except the $2+4$ arithmetic.

\key{6.\ What is computed, and what is proved}

The frames, $m_C$, $q$ and the point $V$ are computed in floating point; the inequality \eqref{eq:p},
checked with an explicit margin (Appendix A of the preprint), is what makes them usable. Two checks
can fail, and both lead to refusal rather than exclusion: if $B_0$ or $C_0$ has no two-dimensional
kernel with room to spare (its rows nearly proportional; for a block with unit entries this is the
$|S|<4$ of the exercise), the frames are not defined; if $q\ge6$,
then $Y$ need not be invertible and \eqref{eq:p} says nothing. Either way the seed is left
\emph{unresolved} --- refined later, never excluded --- and a run that ends with an unresolved seed
has not finished the proof. A failed check is never a proof of anything.

\medskip
\noindent\textit{Exercise.} (a) Check \eqref{eq:blocks} by multiplying out $H'^*H'=6I_6$ in blocks.
(b) With $S=\sum_jb_j$, show that $B'B'^*=\begin{pmatrix}4&\bar S\\ S&4\end{pmatrix}$, so the stretching
factors of $B'$ are $\sqrt{4\pm\abs S}$, and that for an exact $H$ they lie between $\sqrt2$ and
$\sqrt6$. (c) Show that \eqref{eq:p} is sharp: for $\ell_1=q$, $\ell_2=0$ both sides are equal.

\end{document}
